\chapter{Binding（路由绑定）}

\section{概述}

Binding 是将入站消息路由到指定 Agent 的规则。在多智能体模式下，Binding 决定哪个"大脑"处理哪些消息。

配置路径：\texttt{agents.list[].bindings[]}（或顶层 \texttt{bindings[]}）

\section{路由优先级}

路由是\textbf{确定性的}，\textbf{最具体的优先}：

\begin{center}
\begin{tikzpicture}[
    node distance=0.5cm,
    rule/.style={draw=blue!50, fill=blue!5, rounded corners=3pt, minimum width=12cm, minimum height=0.6cm, font=\small, align=left},
    arr/.style={-{Stealth[length=2mm]}, thick, draw=blue!40}
]
    \node[rule] (r1) {1. \texttt{peer} 精确匹配（私信/群组/频道 ID）};
    \node[rule, below=of r1] (r2) {2. \texttt{guildId} 匹配（Discord）};
    \node[rule, below=of r2] (r3) {3. \texttt{teamId} 匹配（Slack）};
    \node[rule, below=of r3] (r4) {4. \texttt{accountId} 匹配（渠道账户）};
    \node[rule, below=of r4] (r5) {5. 渠道级匹配（\texttt{accountId: "*"}）};
    \node[rule, below=of r5] (r6) {6. 默认 Agent（\texttt{agents.list[].default} 或第一个）};

    \draw[arr] (r1) -- (r2);
    \draw[arr] (r2) -- (r3);
    \draw[arr] (r3) -- (r4);
    \draw[arr] (r4) -- (r5);
    \draw[arr] (r5) -- (r6);
\end{tikzpicture}
\end{center}

\section{核心概念}

\begin{center}
\begin{tabular}{ll}
\toprule
\textbf{术语} & \textbf{说明} \\
\midrule
\texttt{agentId} & 一个"大脑"（工作区 + 认证 + 会话） \\
\texttt{accountId} & 渠道账户实例（如 WhatsApp 的 \texttt{"personal"} vs \texttt{"biz"}） \\
\texttt{binding} & 通过 (channel, accountId, peer) 路由到 agentId \\
\texttt{peer} & 对端标识（\texttt{kind: "dm"|"group"|"channel"} + \texttt{id}） \\
\bottomrule
\end{tabular}
\end{center}

\section{配置示例}

\subsection{两个 WhatsApp 账户 → 两个 Agent}

\begin{lstlisting}[language=json]
{
  agents: {
    list: [
      { id: "home", default: true, workspace: "~/.openclaw/workspace-home" },
      { id: "work", workspace: "~/.openclaw/workspace-work" }
    ]
  },
  bindings: [
    { agentId: "home", match: { channel: "whatsapp", accountId: "personal" } },
    { agentId: "work", match: { channel: "whatsapp", accountId: "biz" } }
  ]
}
\end{lstlisting}

\subsection{一个号码按发送者分流}

\begin{lstlisting}[language=json]
{
  agents: {
    list: [
      { id: "alex", workspace: "~/.openclaw/workspace-alex" },
      { id: "mia", workspace: "~/.openclaw/workspace-mia" }
    ]
  },
  bindings: [
    { agentId: "alex", match: { channel: "whatsapp", peer: { kind: "dm", id: "+15551230001" } } },
    { agentId: "mia", match: { channel: "whatsapp", peer: { kind: "dm", id: "+15551230002" } } }
  ]
}
\end{lstlisting}

\subsection{特定群组路由到特定 Agent}

\begin{lstlisting}[language=json]
{
  bindings: [
    {
      agentId: "work",
      match: {
        channel: "whatsapp",
        accountId: "personal",
        peer: { kind: "group", id: "1203630...@g.us" }
      }
    }
  ]
}
\end{lstlisting}

\section{验证}

\begin{lstlisting}[language=bash]
openclaw agents list --bindings
\end{lstlisting}
